\section{Theoretical Background and Related Work}
\label{sec:related}

\paragraph{MMOP test suites and metrics} The MMF suite~\parencite{yue2018moring,yue2019mmf}
and the CEC2019/2020 MMO competitions are standard for MMOP~\parencite{tanabe2020review}. Beyond
objective-space indicators (IGD, IGD+~\parencite{ishibuchi2015igdplus}, hypervolume), MMOP requires
decision-space indicators such as IGDX and the Pareto-set proximity (PSP). Selection in MOEAs
splits broadly into dominance-based (NSGA-II) and decomposition-based
(MOEA/D~\parencite{zhang2007moead}) schemes; EARS is dominance-based, and our contribution is about
\emph{where} a diversity signal enters that dominance pipeline. \textcite{javadi2022interfront}
formalise this distinction for Pareto-based MMOEAs,
classifying selection operations as \emph{inter-front} (acting across fronts, e.g.\ using
neighbours from other fronts to estimate crowding) or \emph{intra-front} (acting only within a
front); in their taxonomy our sparsity signal is an intra-front operation. Our contribution is a
\emph{controlled} comparison of the two placements for a sparsity signal under matched
$E$ and $S$ (Section~\ref{sec:novelty}), realised inside an equivalence-aware framework.

\paragraph{Decision-space niching MOEAs} DN-NSGA-II~\parencite{liang2016dnnsga2} performs
NSGA-II~\parencite{deb2002nsga2} with decision-space crowding; Omni-optimizer~\parencite{deb2008omni}
combines $\epsilon$-dominance with crowding in both spaces; MO\_Ring\_PSO\_SCD
\parencite{yue2018moring} uses a ring-topology PSO with a special crowding distance (SCD) that
fuses objective- and decision-space crowding.

\paragraph{Convergence vs.\ diversity in decision space} CPDEA~\parencite{liu2020cpdea}
introduces a convergence-penalised decision-space density to retain decision-diverse
solutions; MMEA-WI~\parencite{li2021mmeawi} folds decision-space diversity into a weighted
indicator; TriMOEA-TA\&R~\parencite{liu2019trimoea} uses two archives and recombination; and
HREA~\parencite{li2023hrea} ranks hierarchically to handle local Pareto fronts. Recent work
extends MMOP to higher-dimensional decision spaces with dual-population differential
evolution~\parencite{liang2024higherdim}; we examine that regime in Section~\ref{sec:highdim} (where, among the methods compared, EARS's
convergence degrades the least with dimension), though we do not
port that specific dual-population algorithm as a baseline. These methods improve one side of the convergence/diversity
trade-off, typically at a cost to the other (e.g.\ CPDEA attains strong decision-space coverage
but weak objective-space convergence, as we confirm in Section~\ref{sec:experiments}).

\paragraph{Gap} Prior methods let the decision-space-diversity signal enter the
\emph{convergence} decision: CPDEA through a convergence penalty, MMEA-WI through its
selection indicator, and HREA~\parencite{li2023hrea} through a hierarchical ranking that reorders
solutions across local-front structure. EARS-MMOEA instead confines that signal to
\emph{post-sort within-front (intra-front) selection}, leaving the dominance sort untouched; this is the
distinction we make precise in Section~\ref{sec:method}. On the case-study side, MMOP has
seen limited real-network validation~\parencite{tanabe2020review}; we contribute a real-map
facility-placement case study with genuine decision-space multimodality.

We benchmark against three recent state-of-the-art methods (CPDEA, MMEA-WI, and HREA) and three widely-used
classic methods (MO\_Ring\_PSO\_SCD, DN-NSGA-II, Omni-optimizer). Our HREA is a
faithful-in-spirit Python re-implementation of its hierarchy-ranking principle, validated to
behave competently (it recovers multiple equivalent Pareto sets and is a strong competitor,
second overall on the placement case study); like the other re-implemented baselines, any
fidelity gap is disclosed rather than hidden.
